\documentclass[11pt,a4paper]{article}
\usepackage{fontspec,amsmath,amssymb,amsthm,microtype}
\setmainfont{lmroman10-regular.otf}[BoldFont=lmroman10-bold.otf,ItalicFont=lmroman10-italic.otf,BoldItalicFont=lmroman10-bolditalic.otf]
\usepackage[margin=25mm]{geometry}
\usepackage[unicode,hidelinks]{hyperref}
\hypersetup{pdfauthor={LangevinLean Project}}
\setlength{\emergencystretch}{3em}
\setlength{\parskip}{0.4em}
\setlength{\parindent}{0pt}
\allowdisplaybreaks[2]
\raggedbottom
\title{\texorpdfstring{Gaussian and Brownian noise moments}{Gaussian and Brownian noise moments}}
\author{LangevinLean Project}
\date{}
\begin{document}\maketitle
Draft mathematical exposition; these notes may contain mistakes.


Fix a natural-number dimension \(d\in\mathbb N\), write \(E_d=\mathbb R^d\) with its Euclidean norm, and let \(P\) be a measure on \(\Omega\). A law statement \(X_i\sim\mathcal N(0,v)\) below means that the coordinate is almost everywhere measurable and has the stated pushforward measure. We prove the required \(L^2\) norms from these laws. Coordinate independence is unnecessary for second moments because \(\|X\|^2=\sum_iX_i^2\) has no cross terms.



For Brownian processes the input is the scalar Brownian increment laws, whose marginal and increment laws supply the Gaussian hypotheses. These facts establish the noise moments; existence and stationarity of a coupled Langevin reference require separate inputs.



The \(L^2\) norms and expectations below concern the actual measure \(P\), with square-integrability giving square integrability before passing to a quotient class. Laws are pushforward identities; independence is used only for cross-term cancellation, while marginal Gaussian second moments supply the dimension factor.



All second-moment calculations use the Euclidean identity \(\|v\|^2=\sum_{i=1}^d v_i^2\). Only the individual coordinate variances are needed for this identity: independence between coordinates is not required to sum their squares. Independence becomes a separate issue when identifying the joint Gaussian innovation law. For Brownian increments, replacing the starting time by \(a\) changes no second-moment constant because the variance is the elapsed time.


\section{\texorpdfstring{Second moment of a centered Gaussian}{Second moment of a centered Gaussian}}


For nonnegative variance \(v\), the centered Gaussian law has second moment

\[\int_{\mathbb R}x^2\,d\mathcal N(0,v)(x)=v.\]

The known Gaussian variance is \(v\). Expand variance as \(\int(x-\mathbb E[x])^2\), and insert the Gaussian mean \(\mathbb E[x]=0\).



Simplify \(x-0=x\). The variance integral is now exactly the required second moment.



The variance parameter is the second argument of \(\mathcal N(0,v)\), not its standard deviation. Centering therefore makes the variance integral exactly the squared-variable integral. At \(v=0\), the Gaussian is a point mass at zero and the same identity gives zero, so no positive-variance exception is needed.


\section{\texorpdfstring{Gaussian coordinates give a square-\allowbreak{}integrable vector}{Gaussian coordinates give a square-integrable vector}}


Let \(X:\Omega\to E_d\), and suppose every coordinate satisfies \(X_i\sim\mathcal N(0,v)\) under \(P\), with \(v\ge0\). Then \(X\in L^2(P;E_d)\).



In finite-dimensional Euclidean space, square integrability is equivalent to square integrability of each coordinate. Each coordinate belongs to \(L^2\) because its Gaussian law has a finite second moment; transfer that membership along the equality of laws.



A finite sum of integrable nonnegative coordinate squares is integrable. Conversely each coordinate square is bounded by the squared Euclidean norm. These two inequalities explain the equivalence between vector and coordinate \(L^2\) membership in finite dimension. Coordinate independence is not used for this membership statement.


\section{\texorpdfstring{Sum the coordinate moments}{Sum the coordinate moments}}


Under the same coordinate Gaussian laws,

\[\int_\Omega\|X(\omega)\|^2\,dP(\omega)=dv.\]

For each coordinate \(i\), retain the proved square integrability \(X_i\in L^2(P)\). It will justify integrating the finite sum of squared coordinates.



Change variables using the law of \(X_i\), then apply the scalar Gaussian moment:

\[\int_\Omega X_i^2\,dP=\int_{\mathbb R}x^2\,d\mathcal N(0,v)=v.\]

Expand the squared Euclidean norm and interchange the finite sum and integral:

\[\int_\Omega\|X\|^2\,dP=\sum_{i=0}^{d-1}\int_\Omega X_i^2\,dP.\]

Every summand is integrable by its \(L^2\) membership.



Each of the \(d\) summands equals \(v\); hence their sum is \(dv\).



No cross terms occur in the identity for the squared Euclidean norm. Thus possible correlations between distinct coordinates do not change this second moment; their individual variances determine it. The finite sum may be exchanged with the integral because each square has already been proved integrable.


\section{\texorpdfstring{The Gaussian vector as an L\ensuremath{{}^{2}} class}{The Gaussian vector as an L² class}}


Let \([X]\) be the \(L^2\) equivalence class constructed from \(X\), whose coordinates have laws \(\mathcal N(0,v)\). Then

\[\|[X]\|_{L^2}^2=dv.\]

Express the squared \(L^2\) norm as the self-inner-product, then use the definition of the \(L^2\) inner product as an integral.



The chosen representative of \([X]\) equals \(X\) almost everywhere. Therefore the integrand \(\langle[X](\omega),[X](\omega)\rangle\) agrees almost everywhere with \(\|X(\omega)\|^2\), so the two integrals agree.



Substitute the integral identity and the vector second-moment calculation \(\int\|X\|^2\,dP=dv\).



An \(L^2\) class identifies random vectors that agree almost everywhere. The norm is independent of the chosen representative, so the second-moment calculation for the original random vector computes the class norm as well. The inner-product formula makes that identification explicit without selecting pointwise values on exceptional sets.


\section{\texorpdfstring{The square-\allowbreak{}root Gaussian norm}{The square-root Gaussian norm}}


For the same Gaussian coordinate laws,

\[\|[X]\|_{L^2}=\sqrt{dv}.\]

This includes \(d=0\) and \(v=0\).



Replace \(dv\) by \(\|[X]\|^2\). Since norms are nonnegative, \(\sqrt{\|[X]\|^2}=\|[X]\|\).



The argument takes the unique nonnegative square root of \(dv\ge0\). A sign choice is therefore unnecessary, and the formula remains valid at both degenerate endpoints. This is the exact root mean square size used when placing Gaussian vectors inside the common \(L^2\) reference space.


\section{\texorpdfstring{Include the Langevin diffusion factor}{Include the Langevin diffusion factor}}


If every coordinate of \(X\) the stated pushforward law \(\mathcal N(0,v)\), then the Langevin-scaled vector satisfies

\[\|\sqrt2[X]\|_{L^2}=\sqrt{2d}\sqrt v.\]

Norm homogeneity gives \(\|\sqrt2[X]\|=\sqrt2\|[X]\|\), because \(\sqrt2\ge0\). Insert \(\|[X]\|=\sqrt{dv}\), then split the square roots of the nonnegative products.



Rearrange \(\sqrt2\sqrt d\sqrt v\) into \(\sqrt{2d}\sqrt v\).



Multiplying a random vector by \(\sqrt2\) doubles its second moment. Taking a square root changes \(dv\) into \(\sqrt{2dv}\), which is equal to the displayed product because the factors are nonnegative. This factor is the diffusion normalization producing a unit Laplacian in the Langevin generator.


\section{\texorpdfstring{Assemble Brownian coordinates}{Assemble Brownian coordinates}}


For scalar processes \(W_i:[0,\infty)\times\Omega\to\mathbb R\), define the Euclidean vector

\[W_t(\omega):=(W_0(t,\omega),\ldots,W_{d-1}(t,\omega))\in E_d.\]

Collect the coordinates at the fixed time \(t\) and outcome \(\omega\), using the finite-dimensional \(\ell^2\) norm.



Assembling the coordinates changes the codomain from a finite family of scalar variables to \(\mathbb R^d\). Evaluating coordinate \(i\) of the resulting vector simply recovers the original scalar variable at that index. This observation lets every coordinate law and measurability property transfer without a new probabilistic construction.


\section{\texorpdfstring{Brownian marginals are Gaussian}{Brownian marginals are Gaussian}}


Assume each scalar \(W_i\) satisfies \(\text{the Brownian increment laws}\) under \(P\). For each \(t\ge0\) and coordinate \(i\),

\[W_t^{(i)}\sim\mathcal N(0,t).\]

Apply the scalar Brownian marginal-law theorem at time \(t\). The \(i\)-th coordinate of the assembled vector is by definition the original scalar process.



For a pre-Brownian coordinate, the distribution at time \(t\) is already specified by its centered Gaussian marginal of variance \(t\). Selecting the coordinate of the assembled vector therefore gives the claimed law directly. This statement alone makes no assertion about independence between distinct spatial coordinates.


\section{\texorpdfstring{Brownian vectors belong to L\ensuremath{{}^{2}}}{Brownian vectors belong to L²}}


Under the scalar pre-Brownian assumptions, for every \(t\ge0\),

\[W_t\in L^2(P;E_d).\]

Apply Gaussian vector square integrability with variance \(v=t\), using the just-proved marginal laws of all coordinates.



Each coordinate has finite second moment \(t\). Summing the \(d\) coordinate integrals bounds the vector square and gives finite \(L^2\) membership. The conclusion is a fixed-time moment statement; continuity of the process as an \(L^2\) curve is established separately from increment laws.


\section{\texorpdfstring{The exact Brownian noise norm}{The exact Brownian noise norm}}


For scalar pre-Brownian coordinates and \(t\ge0\),

\[\|\sqrt2[W_t]\|_{L^2}=\sqrt{2d}\sqrt t.\]

Specialize the scaled Gaussian vector norm to the Brownian vector, whose coordinate variance is \(t\). This derives the square-root noise scale from the Brownian laws.



The squared norm is exactly \(2dt\), not merely bounded by it. Consequently the reference noise estimate uses the optimal moment scale supplied by the Gaussian law. At \(t=0\), the noise class has zero norm and hence equals zero in \(L^2\), agreeing with the integral-equation initialization.


\section{\texorpdfstring{Define an increment from any starting time}{Define an increment from any starting time}}


For \(a,t\ge0\), define the vector increment over the interval \([a,a+t]\) by

\[\Delta_aW_t(\omega):=W_{a+t}(\omega)-W_a(\omega).\]

Subtract the two Euclidean vectors at the interval endpoints. Each coordinate is the corresponding scalar Brownian increment.



The elapsed time parameter is kept separate from the starting time. This permits every reference step to be described on its local interval \([0,h]\), while its underlying Brownian input is taken from the global interval \([a,a+h]\). Vector subtraction is performed coordinatewise, so scalar increment laws apply immediately.


\section{\texorpdfstring{Increment variance is elapsed time}{Increment variance is elapsed time}}


Under the scalar pre-Brownian assumptions, for every \(a,t\ge0\) and every coordinate \(i\),

\[W_i(a+t)-W_i(a)\sim\mathcal N(0,t).\]

Expand the vector increment to its scalar coordinate. Apply the Brownian difference-law theorem at times \(a+t\) and \(a\); its variance is the nonnegative distance between those times.



It remains to identify that variance. As real numbers,

\[|(a+t)-a|=|t|=t,\]

because \(t\ge0\). Thus the increment law is exactly \(\mathcal N(0,t)\).



The distance of the two nonnegative time points is \(t\), because their difference is \((a+t)-a=t\ge0\). Thus the variance depends on elapsed time alone. This invariance of the increment law under a change of starting time is the probabilistic fact needed to give all reference segments the same error constant.


\section{\texorpdfstring{Increments are square integrable}{Increments are square integrable}}


For pre-Brownian coordinates and any \(a,t\ge0\),

\[\Delta_aW_t\in L^2(P;E_d).\]

Each increment coordinate has the centered Gaussian law of variance \(t\), so the Gaussian vector integrability theorem applies.



Apply the coordinate-to-vector moment criterion to the increment rather than to the two endpoints separately. Each increment coordinate has second moment equal to elapsed time, so the argument gives exactly the moment scale needed later. It does not require a bound growing with the absolute starting time \(a\).


\section{\texorpdfstring{The same noise norm on every reference segment}{The same noise norm on every reference segment}}


For scalar pre-Brownian coordinates and every \(a,t\ge0\),

\[\boxed{\|\sqrt2[W_{a+t}-W_a]\|_{L^2}=\sqrt{2d}\sqrt t.}\]

Apply the scaled Gaussian vector norm to the increment coordinate laws. The answer depends only on elapsed time \(t\), so every reference interval has the same moment bound.



The expression is independent of \(a\), which is crucial when applying the local-defect estimate at step indices \(k=0,1,\ldots\). Each segment has the same \(\sqrt t\) modulus in \(L^2\), so the resulting defect bound is uniform over the entire discrete trajectory. This uniformity supports one fixed forcing constant in the recurrence.


\end{document}
